% !TeX program = pdflatex
\documentclass[11pt,a4paper]{amsart}

\newcommand{\notelanguage}{finnish}
% Shared preamble for course notes and exercise sets.

\usepackage[T1]{fontenc}
\usepackage{lmodern}
\providecommand{\notelanguage}{english}
\usepackage[finnish,english]{babel}
\AtBeginDocument{\expandafter\selectlanguage\expandafter{\notelanguage}}

\newcommand{\notesfinnishlanguage}{finnish}
\ifx\notelanguage\notesfinnishlanguage
\newcommand{\notestheoremname}{Lause}
\newcommand{\noteslemmaname}{Lemma}
\newcommand{\notespropositionname}{Propositio}
\newcommand{\notescorollaryname}{Korollaari}
\newcommand{\notesdefinitionname}{Määritelmä}
\newcommand{\notesexamplename}{Esimerkki}
\newcommand{\notesproblemname}{Tehtävä}
\newcommand{\notesexercisename}{Tehtävä}
\newcommand{\notesremarkname}{Huomautus}
\newcommand{\notessolutionname}{Ratkaisu}
\else
\newcommand{\notestheoremname}{Theorem}
\newcommand{\noteslemmaname}{Lemma}
\newcommand{\notespropositionname}{Proposition}
\newcommand{\notescorollaryname}{Corollary}
\newcommand{\notesdefinitionname}{Definition}
\newcommand{\notesexamplename}{Example}
\newcommand{\notesproblemname}{Problem}
\newcommand{\notesexercisename}{Exercise}
\newcommand{\notesremarkname}{Remark}
\newcommand{\notessolutionname}{Solution}
\fi

\usepackage[a4paper,margin=30mm]{geometry}
\usepackage{microtype}
\usepackage{graphicx}
\usepackage{enumitem}

\newcommand{\exerciseheading}[2]{%
  \par\noindent
  {\large\bfseries #1\par}
  \vspace{0.1em}
  \noindent{\large #2\par}
  \vspace{1.5em}
}

\usepackage{amsmath}
\usepackage{amssymb}
\usepackage{amsthm}
\usepackage{mathtools}
\usepackage{aliascnt}

% Row operations beside the affected row, e.g. \rowop{+2R_1}.
\newcommand{\rowop}[1]{\mathrel{\big|}\,#1}

% Linear systems: the argument gives the number of alignment pairs.
\newenvironment{linearsystem}[1]
{\left\{
  \begin{alignedat}{#1}}
    {
  \end{alignedat}\right.}

  \usepackage[hidelinks,pdfusetitle]{hyperref}

  \numberwithin{equation}{section}
  \setlist{itemsep=0.25em,topsep=0.5em}

  \theoremstyle{plain}
  \newtheorem{theorem}{\notestheoremname}[section]
  \newaliascnt{lemma}{theorem}
  \newtheorem{lemma}[lemma]{\noteslemmaname}
  \aliascntresetthe{lemma}
  \newaliascnt{proposition}{theorem}
  \newtheorem{proposition}[proposition]{\notespropositionname}
  \aliascntresetthe{proposition}
  \newaliascnt{corollary}{theorem}
  \newtheorem{corollary}[corollary]{\notescorollaryname}
  \aliascntresetthe{corollary}

  \theoremstyle{definition}
  \newaliascnt{definition}{theorem}
  \newtheorem{definition}[definition]{\notesdefinitionname}
  \aliascntresetthe{definition}
  \newaliascnt{example}{theorem}
  \newtheorem{example}[example]{\notesexamplename}
  \aliascntresetthe{example}
  \newaliascnt{problem}{theorem}
  \newtheorem{problem}[problem]{\notesproblemname}
  \aliascntresetthe{problem}
  \newtheorem{exercise}{\notesexercisename}

  \theoremstyle{remark}
  \newaliascnt{remark}{theorem}
  \newtheorem{remark}[remark]{\notesremarkname}
  \aliascntresetthe{remark}

  \newenvironment{solution}
  {
  \begin{proof}[\notessolutionname]}
    {
  \end{proof}}

  \usepackage[nameinlink,noabbrev]{cleveref}

  \crefname{theorem}{theorem}{theorems}
  \crefname{lemma}{lemma}{lemmas}
  \crefname{proposition}{proposition}{propositions}
  \crefname{corollary}{corollary}{corollaries}
  \crefname{definition}{definition}{definitions}
  \crefname{example}{example}{examples}
  \crefname{problem}{problem}{problems}
  \crefname{exercise}{exercise}{exercises}
  \crefname{remark}{remark}{remarks}

  \ifx\notelanguage\notesfinnishlanguage
  \crefname{theorem}{lause}{lauseet}
  \crefname{lemma}{lemma}{lemmat}
  \crefname{proposition}{propositio}{propositiot}
  \crefname{corollary}{korollaari}{korollaarit}
  \crefname{definition}{määritelmä}{määritelmät}
  \crefname{example}{esimerkki}{esimerkit}
  \crefname{problem}{tehtävä}{tehtävät}
  \crefname{exercise}{tehtävä}{tehtävät}
  \crefname{remark}{huomautus}{huomautukset}
  \fi

  \newcommand{\NN}{\mathbb{N}}
  \newcommand{\NNone}{\mathbb{N}_1}
  \newcommand{\ZZ}{\mathbb{Z}}
  \newcommand{\QQ}{\mathbb{Q}}
  \newcommand{\RR}{\mathbb{R}}
  \newcommand{\CC}{\mathbb{C}}
  \newcommand{\PP}{\mathbb{P}}

  \DeclareMathOperator{\im}{im}
  \DeclareMathOperator{\rank}{rank}
  \DeclareMathOperator{\syt}{syt}
  \DeclareMathOperator{\pyj}{pyj}
  \DeclareMathOperator{\dom}{dom}
  \DeclarePairedDelimiter{\abs}{\lvert}{\rvert}
  \DeclarePairedDelimiter{\norm}{\lVert}{\rVert}
  \DeclarePairedDelimiterXPP{\maxnorm}[1]{}{\lVert}{\rVert}{_\infty}{#1}
  \newcommand{\set}[1]{\left\lbrace#1\right\rbrace}
  \newcommand{\maxset}[1]{\max\set{#1}}
  \newcommand{\minset}[1]{\min\set{#1}}
  \newcommand{\powerset}[1]{\mathcal{P}\!\left(#1\right)}


\title{Lukuteoria I\\Harjoitus 4}
\author{}
\date{}

\begin{document}

\exerciseheading{Lukuteoria I}{Harjoitus 4}

\begin{exercise}
  \leavevmode\par
  \begin{enumerate}[label=(\alph*), topsep=0.35em]
    \item Määrää rationaaliluvun \(\frac{71}{15}\) desimaaliesitys.
    \item Määrää luvun \(7{,}\overline{12} = 7{,}121212\ldots\)
      esitys murtolukuna.
  \end{enumerate}
\end{exercise}

\begin{solution}
  \leavevmode\par
  \begin{enumerate}[label=(\alph*), topsep=0.35em]
    \item
      \[
        \renewcommand{\arraystretch}{1.05}
        \begin{array}{r@{\;}r@{\mkern3mu}r@{\mkern3mu}r@{\mkern3mu}r@{\mkern3mu}r@{\mkern3mu}l}
          & & 4\mathrlap{{,}} & 7 & 3 & 3 & \ldots \\
          \cline{2-7}
        15\,\big) & 7 & 1\mathrlap{{,}} & 0 & 0 & 0 & \ldots \\
        & 6 & 0 & & & & \\
        \cline{2-3}
        & 1 & 1 & 0 & & & \\
        & 1 & 0 & 5 & & & \\
        \cline{2-4}
        & & & 5 & 0 & & \\
        & & & 4 & 5 & & \\
        \cline{4-5}
        & & & & 5 & 0 & \\
        & & & & 4 & 5 & \\
        \cline{5-6}
        & & & & & 5 & \ldots
      \end{array}
    \]
    \[
      \boxed{\frac{71}{15}=4{,}73333\ldots=4{,}7\overline{3}.}
    \]
  \item Olkoon \(x=7{,}121212\ldots\). Tällöin
    \[
      \begin{aligned}
        100x-x
        &=712{,}121212\ldots-7{,}121212\ldots\\
        &=705.
      \end{aligned}
    \]
    Siis \(99x=705\), joten
    \[
      \boxed{x=\frac{705}{99}=\frac{235}{33}.}\qedhere
    \]
\end{enumerate}
\end{solution}

\begin{exercise}
\leavevmode\par
\begin{enumerate}[label=(\alph*), topsep=0.35em]
  \item Osoita, että \(\sqrt{2} + \sqrt{5}\) on irrationaaliluku.
  \item Osoita, että \(\sqrt[3]{2} + \sqrt{3}\) on algebrallinen luku.
\end{enumerate}
\end{exercise}

\begin{solution}
\leavevmode\par
\begin{enumerate}[label=(\alph*), topsep=0.35em]
  \item Oletetaan, että \(\sqrt{2} + \sqrt{5} = r \in \QQ\). Tällöin
    \((r - \sqrt{2})^2 = 5\), eli
    \[
      r^2 - 2r\sqrt{2} + 2 = 5.
    \]
    Koska \(r>0\), saadaan
    \[
      \sqrt{2} = \frac{r^2 - 3}{2r} \in \QQ.
    \]
    Tämä on ristiriita, sillä \(\sqrt{2}\) on irrationaalinen.
    Siis \(\sqrt{2} + \sqrt{5}\) on irrationaalinen.
  \item Olkoon \(x=\sqrt[3]{2}+\sqrt{3}\). Tällöin
    \((x-\sqrt{3})^3=2\), joten
    \[
      x^3-3x^2\sqrt{3}+9x-3\sqrt{3}=2.
    \]
    Ryhmittelemällä saadaan
    \[
      x^3+9x-2=3(x^2+1)\sqrt{3}.
    \]
    Korottamalla puolittain neliöön saadaan
    \[
      (x^3+9x-2)^2-27(x^2+1)^2=0.
    \]
    Siis \(x\) on kokonaislukukertoimisen polynomin
    \[
      P(T)=(T^3+9T-2)^2-27(T^2+1)^2
    \]
    nollakohta. Täten ollen \(\sqrt[3]{2}+\sqrt{3}\) on
    algebrallinen luku.
\end{enumerate}
\end{solution}

\begin{exercise}
Perustele ovatko seuraavat väitteet tosia vai epätosia.
\begin{enumerate}[label=(\alph*), topsep=0.35em]
  \item Rationaaliluvun ja irrationaaliluvun summa on aina
    irrationaaliluku.
  \item Kahden irrationaaliluvun tulo on aina irrationaaliluku.
  \item Jos \(z\) on transkendenttinen luku, niin myös \(z^2\) on
    transkendenttinen luku.
\end{enumerate}
\end{exercise}

\begin{solution}
\leavevmode\par
\begin{enumerate}[label=(\alph*), topsep=0.35em]
  \item Väite on tosi. Olkoon \(a \in \QQ\) ja
    \(r \in \RR \setminus \QQ\). Oletetaan vastoin väitettä, että
    \(a+r=q\in\QQ\). Tällöin \(r=q-a\in\QQ\), sillä
    rationaalilukujen erotus on rationaalinen. Tämä on ristiriidassa
    oletuksen \(r\notin\QQ\) kanssa. Siis \(a+r\) on irrationaalinen.
  \item Väite on epätosi. Tiedetään, että \(\sqrt{2} \in \RR \setminus \QQ\),
    mutta \(\sqrt{2}\sqrt{2} = 2 \in \QQ\).
  \item Väite on tosi. Olkoon \(z\) transkendenttinen. Oletetaan vastoin
    väitettä, että \(z^2\) on algebrallinen. Tällöin on olemassa
    polynomi \(P \in \ZZ[T]\),
    \[
      P(T) = a_dT^d + \cdots + a_1T + a_0, \quad a_d \neq 0
    \]
    siten, että \(P(z^2) = 0\). Määritellään
    \[
      Q(T) = P(T^2) = a_dT^{2d} + \cdots + a_1T^2 + a_0.
    \]
    Tällöin \(Q \in \ZZ[T]\), ja sen johtokerroin on \(a_d\neq0\),
    joten \(Q\) ei ole nollapolynomi. Lisäksi \(Q(z)=P(z^2)=0\),
    joten \(z\) on algebrallinen. Tämä on ristiriidassa oletuksen
    kanssa, että \(z\) on transkendenttinen. Siis \(z^2\) on
    transkendenttinen.

\end{enumerate}
\end{solution}

\begin{exercise}
Oletetaan tunnetuksi, että luvut \(e\) ja \(\pi\) ovat
transkendenttisia. Osoita, että ainakin toinen luvuista \(e + \pi\)
ja \(e\pi\) on irrationaalinen.
\end{exercise}

\begin{solution}
Oletetaan vastoin väitettä, että luvut \(e+\pi\) ja \(e\pi\) ovat
rationaalisia. Kirjoitetaan
\[
  e+\pi=\frac{a}{b},\qquad e\pi=\frac{c}{d},
  \qquad a,c\in\ZZ,\quad b,d\in\ZZ_+.
\]
Olkoon
\[
  P(T) = T^2 - (e + \pi)T + e\pi = (T - e)(T - \pi).
\]
Kertomalla nimittäjien tulolla saadaan polynomi
\[
  Q(T)=bdP(T)\in\ZZ[T].
\]
Koska \(bd\neq0\), polynomi \(Q\) ei ole nollapolynomi. Lisäksi
\(Q(e)=bdP(e)=0\), joten \(e\) on algebrallinen. Tämä on ristiriidassa
luvun \(e\) transkendenttisuuden kanssa.
Täten ainakin toisen luvuista \(e + \pi\) ja \(e\pi\) on oltava
irrationaalinen.
\end{solution}

\begin{exercise}
Osoita, että Liouvillen luku
\[
  \ell = \sum_{n=1}^{\infty} 10^{-n!}
  = 0{,}1100010000000000000000010000\ldots
\]
on irrationaalinen käyttämättä tietoa siitä, että \(\ell\) on
transkendenttinen.
\end{exercise}

\begin{solution}
Oletetaan, että \(\ell = p/q\), missä \(p,q\in\ZZ\) ja \(q>0\).
Olkoon \(N\in\NNone\). Kerrotaan yhtälö \(\ell=p/q\) luvulla \(q10^{N!}\)
ja siirretään sarjan \(N\) ensimmäistä termiä oikealle puolelle. Saadaan
\[
  R_N := q10^{N!}\sum_{n=N+1}^{\infty} 10^{-n!}
  = p10^{N!} - q\sum_{n=1}^{N} 10^{N!-n!}.
\]
Koska \(N!-n!\) on epänegatiivinen kokonaisluku, kun \(1\le n\le N\),
yhtälön oikea puoli on kokonaisluku. Siis \(R_N\in\ZZ\).

Luvut \(n!-N!\), kun \(n\ge N+1\), muodostavat aidosti kasvavan
kokonaislukujonon, jonka ensimmäinen jäsen on \((N+1)!-N!=N\cdot N!\).
Täydentämällä summaa puuttuvilla positiivisilla termeillä saadaan
geometrinen yläraja
\[
  \begin{aligned}
    0 < R_N &= q\sum_{n=N+1}^{\infty} 10^{-(n!-N!)} \\
    &\le q\sum_{k=N\cdot N!}^{\infty} 10^{-k} \\
    &= q10^{-N\cdot N!}\sum_{j=0}^{\infty} \left(\frac{1}{10}\right)^j \\
    &= q10^{-N\cdot N!}\frac{1}{1-\frac{1}{10}}
    = \frac{10q}{9}10^{-N\cdot N!}.
  \end{aligned}
\]
Koska \(q\) on kiinteä, yläraja lähestyy nollaa, kun \(N\to\infty\).
Valitaan \(N\) niin suureksi, että yläraja on pienempi kuin \(1\).
Tällöin \(0<R_N<1\), mikä on ristiriidassa sen kanssa, että \(R_N\)
on kokonaisluku. Siis \(\ell\) on irrationaalinen.
\end{solution}

\begin{exercise}
Tarkastele todistusta Liouvillen luvun transkendenttisuudelle.
Miten todistusta pitäisi muuttaa, jotta saataisiin osoitettua, että
\[
  3\sum_{i=1}^{\infty} 7^{-i!}
\]
on transkendenttinen?
\end{exercise}

\end{document}
